\chapter{Methodik}
\section{Wissenschaftliche Methode Fallstudienanalyse}

\newcolumntype{A}{>{\raggedright\arraybackslash}p{4.3cm}} % Autor/Jahr
\newcolumntype{S}{>{\raggedright\arraybackslash}p{2.8cm}} % System
\newcolumntype{C}{>{\centering\arraybackslash}p{0.55cm}}  % Konzept-Spalten (alle gleich breit)

\newcommand{\lhdr}[2]{\parbox[c][32mm][c]{#1}{\centering\textbf{#2}}}
\newcommand{\hdr}[1]{\rotatebox{90}{\scriptsize\strut #1}}

\newcommand{\yes}{\textbf{x}}

\begin{table}[htbp]
\centering
\scriptsize
\setlength{\tabcolsep}{2.5pt}
\renewcommand{\arraystretch}{1.25}

\begin{tabular}{|A|S||*{14}{C|}}
\hline
\multicolumn{2}{|c||}{\textbf{Studie}} &
\multicolumn{14}{c|}{\textbf{Konzepte}} \\
\hline

\multirow{2}{*}{\lhdr{4.3cm}{Autor(en)\\\& Jahr}} &
\multirow{2}{*}{\lhdr{2.8cm}{System}} &
\multicolumn{3}{c|}{\textbf{Ziel}} &
\multicolumn{3}{c|}{\textbf{Datenquellen}} &
\multicolumn{3}{c|}{\textbf{KI-/LLM-Ansatz}} &
\multicolumn{3}{c|}{\textbf{Evaluation}} &
\multicolumn{2}{c|}{\textbf{Kontext}} \\
\cline{3-16}

& &
\hdr{WCAG-Konf.\ prüfen} &
\hdr{Barrieren korrigieren} &
\hdr{Entwickler unterstützen} &
\hdr{Tool-Reports (axe etc.)} &
\hdr{Interaktionsdaten (Playwright)} &
\hdr{Quellcode-Analyse} &
\hdr{LLM-basiert} &
\hdr{Kontextsensitiv / RAG} &
\hdr{Lokal ausführbar} &
\hdr{Automatisierte Eval.} &
\hdr{Qualitative Eval.} &
\hdr{Nutzerstudie} &
\hdr{Rich-Client / SPA} &
\hdr{Rechtl.\ Rahmen (BFSG etc.)} \\
\hline

% ── HOCH-priorisierte Studien ────────────────────────────────
\multicolumn{16}{|l|}{\textbf{Hoch priorisierte Studien}} \\
\hline

He et al.\ (2025) & GenA11y
& \yes &   &   & \yes &   &   & \yes &   &   & \yes &   &   &   &   \\ \hline

Fathallah et al.\ (2025) & AccessGuru
& \yes & \yes & \yes & \yes & \yes &   & \yes &   &   & \yes &   &   &   &   \\ \hline

Bassi et al.\ (2025) & LLM Auditing
& \yes & \yes & \yes &   &   &   & \yes &   &   &   & \yes &   &   & \yes \\ \hline

Paternò et al.\ (2025) & TeVal
& \yes &   & \yes & \yes &   &   & \yes & \yes &   &   & \yes &   &   &   \\ \hline

Tafreshipour et al.\ (2024) & Ma11y
& \yes &   &   & \yes & \yes &   &   &   &   & \yes &   &   &   &   \\ \hline

% ── Eigener PoC ──────────────────────────────────────────────

\textbf{Martinez Campo (2026)} &
\textbf{CLARA}
& \yes & \yes & \yes
& \yes & \yes & \yes
& \yes & \yes & \yes
&   & \yes & 
& \yes & \yes \\ \hline

\end{tabular}

\caption[Konzeptmatrix (Hauptmatrix, hoch priorisierte Studien)]{%
  Konzeptmatrix nach Webster Watson\footcite[Vgl.][]{WebsterWatson2002}: Hoch priorisierte Studien im Vergleich
  zum eigenen Proof of Concept. \yes~=~Konzept vorhanden/adressiert; leer ~=~nicht adressiert.
  Die vollständige Matrix aller analysierten Studien findet sich in Anhang~\ref{app:konzeptmatrix}.}
\label{tab:konzeptmatrix_haupt}
\end{table}